\chapter{Embedded Firmware and Software}\label{ch:embedded-firmware-and-software}

\section{Repository Organization}\label{repository-organization}

The repository separates portable common logic from platform-specific
implementations under 04\_firmware \cite{zephyr_repo}. The later architecture
includes common/, \path{platforms/esp32}/, and \path{platforms/stm32/}. The common
layer contains interfaces, device-independent drivers, estimation and
control types, and application state; the platform layer owns MCU
peripherals, RTOS tasks, pin configuration, and build files. This
structure supports the stated goal of bringing up the robot on an ESP32
while preserving a path to the STM32H723ZG.

Earlier STM32 development material also contains CubeMX-generated Core,
Drivers, and Middlewares trees, CMake and linker support, and a
zephyr.ioc configuration. Those files represent a historical or
scaffolded platform and should not be described as the active robot
firmware unless the final repository establishes that status. The final
thesis should name the exact target and commit used for every
experiment.

\begin{longtable}[]{@{}
  >{\raggedright\arraybackslash}p{(\columnwidth - 0\tabcolsep) * \real{1.0000}}@{}}
\toprule\noalign{}
\begin{minipage}[b]{\linewidth}\raggedright
\textless INSERT FIGURE 8.1: Create a directory tree for the final
tested commit. Distinguish portable common code, ESP32 BSP and tasks,
STM32 scaffold or implementation, configuration headers, tests\slash{}bring-up
code, documentation, and generated vendor code. Mark historical
directories\textgreater{}\\
What the figure should show: Create a directory tree for the final
tested commit. Distinguish portable common code, ESP32 BSP and tasks,
STM32 scaffold or implementation, configuration headers, tests\slash{}bring-up
code, documentation, and generated vendor code. Mark historical
directories\\
Likely source: repository tree at the final experiment commit\strut
\end{minipage} \\
\midrule\noalign{}
\endhead
\bottomrule\noalign{}
\endlastfoot
\end{longtable}

\begin{center}\singlespacing
\phantomsection\addcontentsline{lof}{figure}{\protect\numberline{8.1}Firmware directory and portability layers.}
\textbf{Figure 8.1. Firmware directory and portability layers.}
\end{center}

\section{Sensor Acquisition and Estimation Path}\label{sensor-acquisition-and-estimation-path}

The portable IMU interface uses \path{common/include/interfaces/if_spi.h},
\path{common/include/drivers/drv_ism330dhcx.h}, and
\path{common/src/drivers/drv_ism330dhcx.c} in the previously inspected
structure. The driver reads the ISM330DHCX identity register and
configures accelerometer and gyroscope control registers. A reported
configuration used 208 Hz output data rate, \ensuremath{\pm}4 g acceleration range,
\ensuremath{\pm}500 degrees\slash{}s angular-rate range, and a 14-byte burst beginning at the
temperature\slash{}gyro\slash{}accelerometer output block. These values must be
rechecked against the final source and the ISM330DHCX data sheet
\cite{st2026ism330dhcx}.

The higher-level path is SPI BSP \ensuremath{\rightarrow} portable IMU driver \ensuremath{\rightarrow} mounting
transform \ensuremath{\rightarrow} gyro calibration \ensuremath{\rightarrow} attitude estimator \ensuremath{\rightarrow} timestamped IMU
snapshot \ensuremath{\rightarrow} safety check \ensuremath{\rightarrow} balance-state extraction. The snapshot should
include a monotonically increasing timestamp, validity flag, sequence
counter or update marker, and diagnostic status. A consumer should copy
one consistent snapshot rather than read fields that can change during
the calculation.

\section{Radio, Display, and User Interaction}\label{radio-display-and-user-interaction}

The RadioMaster RP2 V2 receiver uses the CRSF serial protocol. The
current firmware reportedly decodes 16 channels and monitors link state.
CRSF framing and channel definitions are available in the ExpressLRS
source \cite{expresslrs2026crsf}, while the installed receiver manual and configuration
should be archived with the project \cite{radiomaster2026manuals}. The production command
path should map only documented channels, apply range checks and
deadbands, and force a neutral command when a valid frame is not
received within the selected timeout.

The ST7789 display provides local bring-up and operating status
\cite{waveshare2026lcd}. UI work should run at lower priority than sensing, estimation,
control, and motor communication. It should consume published state
rather than directly reading hardware owned by another task. Diagnostic
pages should show mode, arm state, pitch, IMU status, radio status,
battery state, four actuator states, CAN faults, and timing without
allowing the display driver to block a critical loop.

\section{CAN and Actuator Path}\label{can-and-actuator-path}

The ESP32 platform uses TWAI at 1 Mbit\slash{}s and communicates with CubeMars
actuators. The firmware includes a four-slot actuator snapshot and
feedback\slash{}fault monitoring. The inspected configuration identified motor
1 with CAN ID 69 while the remaining IDs or enable flags were not
confirmed. Servo-velocity commands were issued on a reported 20 ms
period in config\_actuator.h. That interval is evidence of the current
bring-up configuration, not necessarily the final balance-loop rate.

The motor path should have one command owner. The controller publishes
desired physical commands, safety publishes torque permission and fault
state, and the motor task transmits the approved result. Bring-up tests
that directly transmit actuator frames should be compiled or enabled
only in an explicit test build. Feedback parsing should reject invalid
lengths and identifiers, apply mode-specific scaling, timestamp every
update, and expose motor faults without clearing them silently.

\begin{longtable}[]{@{}
  >{\raggedright\arraybackslash}p{(\columnwidth - 0\tabcolsep) * \real{1.0000}}@{}}
\toprule\noalign{}
\begin{minipage}[b]{\linewidth}\raggedright
\textless INSERT FIGURE 8.2: Show producers and consumers of IMU
samples, estimated state, CRSF command, supervisor mode, safety
permission, desired wrench, allocated wheel commands, CAN frames,
actuator feedback, UI state, and logging. Identify snapshot, queue,
notification, or direct-call boundaries\textgreater{}\\
What the figure should show: Show producers and consumers of IMU
samples, estimated state, CRSF command, supervisor mode, safety
permission, desired wrench, allocated wheel commands, CAN frames,
actuator feedback, UI state, and logging. Identify snapshot, queue,
notification, or direct-call boundaries\\
Likely source: final firmware architecture and source review\strut
\end{minipage} \\
\midrule\noalign{}
\endhead
\bottomrule\noalign{}
\endlastfoot
\end{longtable}

\begin{center}\singlespacing
\phantomsection\addcontentsline{lof}{figure}{\protect\numberline{8.2}Runtime data flow through sensing, estimation, supervision, control, and actuation.}
\textbf{Figure 8.2. Runtime data flow through sensing, estimation, supervision, control, and actuation.}
\end{center}

\section{Initialization and Runtime Scheduling}\label{initialization-and-runtime-scheduling}

A safe initialization order begins with outputs in inactive states,
configuration loading, timebase and logging initialization,
communication peripherals, sensor identity checks, actuator bus startup,
calibration, and task creation. The supervisor remains disarmed until
the IMU estimate is valid, actuator feedback is fresh, the operator link
is valid, no safety fault is active, and pitch is within the arm window.
Startup should never preserve a prior nonzero command.

\begin{longtable}[]{@{}
  >{\raggedright\arraybackslash}p{(\columnwidth - 0\tabcolsep) * \real{1.0000}}@{}}
\toprule\noalign{}
\begin{minipage}[b]{\linewidth}\raggedright
\textless INSERT FIGURE 8.3: Draw the exact initialization sequence from
reset to DISARMED and READY. Include GPIO safe levels, SPI\slash{}IMU identity,
gyro calibration, TWAI start, motor discovery, radio start, display
start, configuration checks, snapshot validity, and arm
conditions\textgreater{}\\
What the figure should show: Draw the exact initialization sequence from
reset to DISARMED and READY. Include GPIO safe levels, SPI\slash{}IMU identity,
gyro calibration, TWAI start, motor discovery, radio start, display
start, configuration checks, snapshot validity, and arm conditions\\
Likely source: final main entry point, task creation, and supervisor
code\strut
\end{minipage} \\
\midrule\noalign{}
\endhead
\bottomrule\noalign{}
\endlastfoot
\end{longtable}

\begin{center}\singlespacing
\phantomsection\addcontentsline{lof}{figure}{\protect\numberline{8.3}Firmware initialization and arming sequence.}
\textbf{Figure 8.3. Firmware initialization and arming sequence.}
\end{center}

\clearpage
\begin{landscape}
\par\medskip\noindent\singlespacing
\phantomsection\addcontentsline{lot}{table}{\protect\numberline{8.1}Runtime tasks and timing requirements}
\textbf{Table 8.1. Runtime tasks and timing requirements}\par

\begin{longtable}[]{@{}
  >{\raggedright\arraybackslash}p{(\columnwidth - 8\tabcolsep) * \real{0.2230}}
  >{\raggedright\arraybackslash}p{(\columnwidth - 8\tabcolsep) * \real{0.2374}}
  >{\raggedright\arraybackslash}p{(\columnwidth - 8\tabcolsep) * \real{0.1942}}
  >{\raggedright\arraybackslash}p{(\columnwidth - 8\tabcolsep) * \real{0.1799}}
  >{\raggedright\arraybackslash}p{(\columnwidth - 8\tabcolsep) * \real{0.1655}}@{}}
\toprule\noalign{}
\begin{minipage}[b]{\linewidth}\raggedright
\textbf{Function\slash{}task}
\end{minipage} & \begin{minipage}[b]{\linewidth}\raggedright
\textbf{Trigger or period}
\end{minipage} & \begin{minipage}[b]{\linewidth}\raggedright
\textbf{Priority\slash{}deadline}
\end{minipage} & \begin{minipage}[b]{\linewidth}\raggedright
\textbf{Worst-case time}
\end{minipage} & \begin{minipage}[b]{\linewidth}\raggedright
\textbf{Status}
\end{minipage} \\
\midrule\noalign{}
\endhead
\bottomrule\noalign{}
\endlastfoot
IMU acquisition and estimation & \textless VERIFY; IMU configured ODR
and task\slash{}event trigger\textgreater{} & Highest sensing\slash{}control class &
\textless MEASURE\textgreater{} & Implemented infrastructure \\
Balance control & \textless SELECT AND MEASURE\textgreater{} & Highest
control class & \textless MEASURE\textgreater{} & Planned \\
Motor command\slash{}feedback & Reported 20 ms servo command period; final
pending & High & \textless MEASURE\textgreater{} & Partial \\
CRSF receive\slash{}decode & UART\slash{}event-driven or configured poll & Medium-high
& \textless MEASURE\textgreater{} & Implemented \\
Supervisor\slash{}safety & \textless VERIFY\textgreater{} & High enough to
enforce timeout & \textless MEASURE\textgreater{} & Infrastructure
reported \\
Display\slash{}UI & \textless VERIFY; target 10--20 Hz if
retained\textgreater{} & Low & \textless MEASURE\textgreater{} &
Implemented \\
Telemetry\slash{}logging & Buffered\slash{}event-driven & Low & \textless MEASURE
including backpressure\textgreater{} & Incomplete or unverified \\
\end{longtable}
\end{landscape}
\clearpage

\emph{Note.} Earlier STM32 prototypes used task periods such as 2 ms
CRSF, 10 ms IMU, 50 ms OLED, and 200 ms application UI. Those historical
values are not attributed to the current ESP32 build without source
confirmation.

\begin{longtable}[]{@{}
  >{\raggedright\arraybackslash}p{(\columnwidth - 0\tabcolsep) * \real{1.0000}}@{}}
\toprule\noalign{}
\begin{minipage}[b]{\linewidth}\raggedright
\textless INSERT FIGURE 8.4: Create a measured timing diagram over
several control cycles showing IMU data-ready or read, estimator
completion, control computation, allocation, CAN transmit, actuator
feedback, radio processing, and low-priority UI\slash{}logging. Include
minimum, mean, maximum, and jitter\textgreater{}\\
What the figure should show: Create a measured timing diagram over
several control cycles showing IMU data-ready or read, estimator
completion, control computation, allocation, CAN transmit, actuator
feedback, radio processing, and low-priority UI\slash{}logging. Include
minimum, mean, maximum, and jitter\\
Likely source: firmware timestamps, logic analyzer, and CAN trace\strut
\end{minipage} \\
\midrule\noalign{}
\endhead
\bottomrule\noalign{}
\endlastfoot
\end{longtable}

\begin{center}\singlespacing
\phantomsection\addcontentsline{lof}{figure}{\protect\numberline{8.4}Control-cycle timing and communication schedule.}
\textbf{Figure 8.4. Control-cycle timing and communication schedule.}
\end{center}

\section{Safety, Fault Handling, and Data Freshness}\label{safety-fault-handling-and-data-freshness}

Safety is implemented as explicit state and permission, not as scattered
limit checks. Each data product carries a timestamp and validity. The
safety logic compares age against a source-specific limit, checks
numeric range and finiteness, and latches faults that require operator
inspection. Motor feedback faults, bus-off, IMU failure, radio loss, low
battery, excessive pitch, and emergency stop all remove actuator
permission.

\par\medskip\noindent\singlespacing
\phantomsection\addcontentsline{lot}{table}{\protect\numberline{8.2}Proposed operating and safety states}
\textbf{Table 8.2. Proposed operating and safety states}\par

\begin{longtable}[]{@{}
  >{\raggedright\arraybackslash}p{(\columnwidth - 6\tabcolsep) * \real{0.1460}}
  >{\raggedright\arraybackslash}p{(\columnwidth - 6\tabcolsep) * \real{0.3066}}
  >{\raggedright\arraybackslash}p{(\columnwidth - 6\tabcolsep) * \real{0.2555}}
  >{\raggedright\arraybackslash}p{(\columnwidth - 6\tabcolsep) * \real{0.2920}}@{}}
\toprule\noalign{}
\begin{minipage}[b]{\linewidth}\raggedright
\textbf{State}
\end{minipage} & \begin{minipage}[b]{\linewidth}\raggedright
\textbf{Entry condition}
\end{minipage} & \begin{minipage}[b]{\linewidth}\raggedright
\textbf{Permitted actuator behavior}
\end{minipage} & \begin{minipage}[b]{\linewidth}\raggedright
\textbf{Exit condition}
\end{minipage} \\
\midrule\noalign{}
\endhead
\bottomrule\noalign{}
\endlastfoot
DISARMED & Reset, user disarm, or safe recovery & Torque disabled or
documented zero-output mode & Valid system and explicit prepare\slash{}arm
sequence \\
INITIALIZING & Peripheral and estimator startup & Disabled & Calibration
and device checks complete or fault \\
READY & Valid sensors and motors; within arm window & Disabled &
Explicit arm request or new fault \\
BALANCING & Arm request with all guards true & Safety-approved
controller command & Disarm, fall, stale data, or fault \\
FALLEN & Pitch\slash{}rate exceeds recovery boundary & Disabled\slash{}zero & Manual
reset after safe repositioning \\
FAULT & Electrical, communication, sensor, or software fault &
Disabled\slash{}zero & Fault cleared and explicit manual reset \\
\end{longtable}

\begin{longtable}[]{@{}
  >{\raggedright\arraybackslash}p{(\columnwidth - 0\tabcolsep) * \real{1.0000}}@{}}
\toprule\noalign{}
\begin{minipage}[b]{\linewidth}\raggedright
\textless INSERT FIGURE 8.5: Trace one desired motor command through
mode checks, data-age checks, pitch limits, E-stop, actuator limits,
packing, TWAI transmission, feedback timeout, and fault transition.
Identify the single function allowed to submit a production CAN
command\textgreater{}\\
What the figure should show: Trace one desired motor command through
mode checks, data-age checks, pitch limits, E-stop, actuator limits,
packing, TWAI transmission, feedback timeout, and fault transition.
Identify the single function allowed to submit a production CAN
command\\
Likely source: final motor task and safety code\strut
\end{minipage} \\
\midrule\noalign{}
\endhead
\bottomrule\noalign{}
\endlastfoot
\end{longtable}

\begin{center}\singlespacing
\phantomsection\addcontentsline{lof}{figure}{\protect\numberline{8.5}Motor-command safety gate and fault response.}
\textbf{Figure 8.5. Motor-command safety gate and fault response.}
\end{center}

\section{Configuration, Testing, and Version Control}\label{configuration-testing-and-version-control}

Hardware-dependent values should be centralized in configuration headers
or generated platform configuration: pin assignments, CAN bitrate and
identifiers, IMU ranges and sample rate, receiver UART settings, display
pins, actuator mode, command period, limits, and safety timeouts. Values
should use physical units at module boundaries. Compile-time assertions
and startup logs can detect invalid combinations.

Production and bring-up paths should remain separate. A bring-up build
may contain direct tests for debug UART, CRSF, IMU identity, display
pages, CAN loopback, and one-motor motion. The normal build should start
through the supervisor and safety path. Each thesis experiment should
record the Git commit, build target, configuration checksum or file
revision, actuator firmware\slash{}settings, and any local changes.

\par\medskip\noindent\singlespacing
\phantomsection\addcontentsline{lot}{table}{\protect\numberline{8.3}Firmware modules and repository evidence}
\textbf{Table 8.3. Firmware modules and repository evidence}\par

\begin{longtable}[]{@{}
  >{\raggedright\arraybackslash}p{(\columnwidth - 6\tabcolsep) * \real{0.3066}}
  >{\raggedright\arraybackslash}p{(\columnwidth - 6\tabcolsep) * \real{0.2482}}
  >{\raggedright\arraybackslash}p{(\columnwidth - 6\tabcolsep) * \real{0.2336}}
  >{\raggedright\arraybackslash}p{(\columnwidth - 6\tabcolsep) * \real{0.2117}}@{}}
\toprule\noalign{}
\begin{minipage}[b]{\linewidth}\raggedright
\textbf{Repository path or module}
\end{minipage} & \begin{minipage}[b]{\linewidth}\raggedright
\textbf{Responsibility}
\end{minipage} & \begin{minipage}[b]{\linewidth}\raggedright
\textbf{Evidence at inspected snapshot}
\end{minipage} & \begin{minipage}[b]{\linewidth}\raggedright
\textbf{Thesis status}
\end{minipage} \\
\midrule\noalign{}
\endhead
\bottomrule\noalign{}
\endlastfoot
\path{04_firmware/common/include/interfaces/if_spi.h} & Portable SPI
abstraction & Previously inspected & Implemented structure \\
\path{04_firmware/common/include/drivers/drv_ism330dhcx.h} & IMU public
interface & Previously inspected & Implemented \\
\path{04_firmware/common/src/drivers/drv_ism330dhcx.c} & ISM330DHCX register
and conversion logic & Previously inspected & Implemented\slash{}verify final
configuration \\
\path{04_firmware/common/include/control/ctrl_balance_types.h} &
Balance-state\slash{}control data types & Observed at commit 97e5a0a &
Interface only; controller absent \\
\path{04_firmware/platforms/esp32} & Active BSP, tasks, and peripheral
integration & Reported working bring-up platform & Partial system \\
\path{04_firmware/platforms/esp32/main/config/config_actuator.h} & Actuator
configuration and command period & 20 ms servo-velocity period
previously observed & Bring-up setting; final value pending \\
\path{04_firmware/platforms/stm32} and\slash{}or generated STM32 tree & Future STM32
implementation & Scaffold\slash{}generated code & Planned migration \\
\path{03_electrical/Zephyr.kicad_*} & Controller PCB design & Schematic\slash{}PCB
project present & Design in progress \\
05\_simulation & Current repository simulation directory & Visible at
repository root; contents not fully re-inspected & Requires ZIP\slash{}current
connector verification \\
\end{longtable}

\begin{longtable}[]{@{}
  >{\raggedright\arraybackslash}p{(\columnwidth - 0\tabcolsep) * \real{1.0000}}@{}}
\toprule\noalign{}
\begin{minipage}[b]{\linewidth}\raggedright
\textless VERIFY FROM REPOSITORY: Replace the summarized paths in Table
8.1 with the exact final tree, task names, line-referenced configuration
values, and tested commit. The current GitHub directory traversal was
not available to the document-generation environment even after
connection; upload a ZIP containing .git metadata or git log output if
branch and commit-history evidence is required.\textgreater{}
\end{minipage} \\
\midrule\noalign{}
\endhead
\bottomrule\noalign{}
\endlastfoot
\end{longtable}
